\documentclass[11pt]{article}
\usepackage[T1]{fontenc}
\usepackage{lmodern,microtype}
\usepackage{amsmath,amssymb,amsthm}
\usepackage[letterpaper,margin=1in]{geometry}
\usepackage{booktabs,enumitem}
\usepackage[colorlinks=true,linkcolor=blue!50!black,urlcolor=blue!50!black]{hyperref}
\newcommand{\tH}{\widetilde H}
\theoremstyle{plain}
\newtheorem*{prop*}{Proposed Corollary}
\title{Referee report on the rank-two toral positivity note\\
\normalsize (with Part~IV, Draft~5, and its revision log)}
\author{Claude (Anthropic), for David Joyner}
\date{November 7, 2026}
\begin{document}
\maketitle

\section*{0. What was submitted}
Three files were submitted.
\begin{itemize}
\item \texttt{mjl\_part4\_draft5.tex}. This is byte-for-byte identical to
the Draft~5 I refereed earlier today (same MD5 checksum). My earlier
report on it stands. In particular, delete the nongood-class hypothesis
in Proposition~8.5, correct the description of $B_d$, and label
Lemma~6.4.
\item The Markdown revision log for Draft~5. It is accurate. The Python check script that it mentions was not attached.
\item The new note, \emph{A permutation-character theorem for quadratic
wild-place discriminant ratios}. This is the substance of this
report.
\end{itemize}
I read all three files in full.

\section*{1. The note: verdict}
\textbf{The theorem is correct.} I checked every step by hand (\S1.1).
I also wrote an independent script that recomputes the coefficients
$c_k$ from the valuation data. It covers $p\in\{3,5,7\}$, $p^a\in\{p,p^2\}$,
several ramification indices and residue degrees for $F$, and both
types of $E$: 72 cases in all. In every case every $c_k$ is a
nonnegative integer, and the telescoped sum equals
$R_E(1)=|n|_F^{-1}$. The examples are also right: $1+104\,\mathrm{Reg}(C_6)$
and $1+12\,\mathrm{Reg}(C_2)+60\,\mathrm{Reg}(C_{10})$ for
$\mathbb Q_5(\zeta_5)$ both have value $625=|5|_F^{-1}$ at the identity.

This settles Step~T of Part~IV, \S6.4, for $r=2$, at every wild place
and for every odd $n$. Together with Lemma~6.4 it implies the
following. For every ordinary irreducible $\rho$ of a quaternion
division algebra, the restriction of $L_D\Theta_\rho$ to every covered
maximal torus is an actual genuine character.

\subsection*{1.1 Checks}
\begin{itemize}
\item \emph{The discriminant.} $\Delta(t)=|N(1-u)|_F^{1/2}=|1-u|_E^{1/2}$,
since $\mathrm{Ad}(t)$ acts on $D/E=Ej$ by multiplication by
$u=t/\sigma(t)$. Then $1-u^n=\prod_{\zeta\in\mu_n}(1-\zeta u)$.
\item \emph{Roots of unity with nontrivial prime-to-$p$ part.} For these,
$|1-\zeta u|_E=1$. The reduction of $u\in E^1$ meets $k_F^\times$ only
in $\{\pm1\}$: in the unramified case because the norm-one group of
$k_E^\times$ meets $k_F^\times$ in $\{\pm1\}$, and in the ramified case
because $\bar u^2=1$. A nontrivial root of unity of odd order is not
$\pm1$.
\item \emph{The radial formula.} $v_E(1-\zeta u)=\min(v_E(1-u),t_\zeta)$. The equality case is
excluded by applying $\sigma$, using $\sigma(u)=u^{-1}$, and
$v_E(\zeta-\zeta^{-1})=t_{\zeta^2}=t_\zeta$, which holds for $p$ odd.
\item \emph{Value at the identity.} $\prod_{\zeta\in\mu_{p^a}\smallsetminus1}(1-\zeta)=p^a$
gives $R_E(1)=|n|_F^{-1}$. This agrees with Part~III, Lemma~3.3, for
$r=2$.
\item \emph{Filtration indices.} Unramified: $[G:G_k]=(q+1)q^{k-1}$.
Ramified: $[G:G_1]=2$ and $[G:G_{2j}]=2q^j$, with all positive depths
odd. The ramified case uses the parametrization $c\mapsto u$, and
$u-u'$ has $E$-valuation $2v_F(c-c')+1$.
\item \emph{Integrality of $c_k$.} It uses three facts:
\begin{itemize}
\item the exponents $s_k$ are even, so $q+1$ divides $q^{s_k}-1$;
\item $q^{k-1}$ divides $M_{k-1}$, respectively $q^j$ divides $M_{2j-1}$;
\item $M_1$ is odd.
\end{itemize}
All correct.
\end{itemize}

\subsection*{1.2 Minor repairs to the note}
\begin{enumerate}[label=(\arabic*)]
\item \emph{Surjectivity at each level.} The unramified filtration
indices are stated without proof that the norm-one group surjects onto
the trace-zero residue groups at every level. Add one line, via Hensel's
lemma, or via the Hilbert~90 parametrization as in the ramified case.
\item \emph{Where $p$ odd is used.} The theorem uses $p$ odd in two
places: $\zeta\ne\zeta^{-1}$, and the $\{\pm1\}$ argument. Say so in the
statement.
\item \emph{Labels.} Give the theorem a \textbf{[P]} label, as in the
series.
\item \emph{Integration.} Make it Proposition~6.6 of Part~IV, replacing
the example-level Proposition~6.5. Then state the consequence for
\emph{all} $\rho$ via Lemma~6.4.
\item \emph{Its last remark.} It is correct: positivity of one row does
not give integrality of that row. This exposes an overstatement in my
Draft~5 report (\S2).
\end{enumerate}

\section*{2. A correction to my previous report}
In my Draft~5 report, \S2(ii), I said the positive-semidefinite test for
the single row $\rho=\mathbf1$ would settle that case. That is too
strong.
\begin{itemize}
\item Positive-semidefiniteness of one row shows only that the row is
nonnegative.
\item The positivity lemma needs \emph{all} rows to be nonnegative,
because disjointness of supports comes from comparing rows.
\end{itemize}
The finite test is still useful: a negative eigenvalue would disprove
integrality at once, and it is cheap. But to \emph{prove}
irreducibility of $L_D\mathbf1$ one needs either integrality of that
row, or nonnegativity of every row whose support meets that of
$L_D\mathbf1$. The lemma itself, (a)$\Leftrightarrow$(c) for the whole
matrix, is unaffected.

\section*{3. Can the wild integrality problem now be solved in general?}
\textbf{Not yet. But it can now be stated much more precisely.} Here is
where things stand.

\paragraph{What is proved.}
\begin{itemize}
\item The unitary transform and its matrix (Part~III).
\item Integrality $\Leftrightarrow$ nonnegativity of the whole matrix
(the positivity lemma).
\item Torus-restriction positivity for every $\rho$ when $r=2$ (this
note together with Lemma~6.4).
\end{itemize}

\paragraph{What is not proved.}
\begin{enumerate}[label=(\alph*)]
\item \emph{Step~T for $r\ge3$.} The note's argument depends on
$\sigma(u)=u^{-1}$ for norm-one elements of a quadratic extension. For
degree $r\ge3$ the root ratios $u_{ij}=\lambda_i/\lambda_j$ lie in the
Galois closure, and non-Galois $E$ occur. The equality case of the
radial formula and the filtration indices both need new arguments.
\item \emph{Step~G, gluing, for every $r$ including $r=2$.} This is the
real obstruction. Positivity on every torus does not imply positivity on
the full group. Positive-definiteness is not determined by restrictions
to abelian subgroups.
\end{enumerate}

\subsection*{3.1 What torus positivity does give: rationality, via Artin}
The following is a short consequence that sharpens what Step~G must
supply. I believe it is correct, but it has not been checked by anyone
else.

\begin{prop*}[\textbf{[P]}, given Step~T on all maximal tori; in
particular for $r=2$]
Every coefficient $a_{\rho\tau}$ is rational, with denominator dividing
$|Q|$ for any finite twisted quotient $Q=\tH/N$ through which
$L_D\Theta_\rho$ and $\Theta_\tau$ factor.
\end{prop*}

\begin{proof}[Proof sketch]
\begin{enumerate}[label=(\roman*)]
\item \emph{Cyclic subgroups lie in covered tori.} Every cyclic subgroup
of $Q$ is the image of some $\langle g\rangle$ with $g\in\tH$. Since $r$
is prime, $p(g)$ lies in a maximal subfield $E$, so
$\langle g\rangle\subset\widetilde{E^\times}=p^{-1}(E^\times)$.
\item \emph{The restriction to a torus is a character of its image.}
By Step~T, $\varphi=L_D\Theta_\rho$ restricts to an actual character of
$\widetilde{E^\times}$. This character is invariant under translation
by $K=\widetilde{E^\times}\cap N$. For a nonnegative combination of
characters, $\varphi(k)=\varphi(1)$ forces $k$ to act trivially in
every constituent. So $\varphi$ restricts to a genuine character of the
image of $\widetilde{E^\times}$ in $Q$, hence of every cyclic subgroup.
\item \emph{Artin's induction theorem.} It gives
$|Q|\cdot1_Q=\sum_Ca_C\,\mathrm{Ind}_C^Q\theta_C$, with $a_C\in\mathbb Z$
and $\theta_C$ ordinary virtual characters. Multiplying by the genuine
function $\varphi$ and using
$\varphi\cdot\mathrm{Ind}_C\theta=\mathrm{Ind}_C(\varphi|_C\,\theta)$
gives $|Q|\varphi\in R^{\mathrm{gen}}(Q)$.
\end{enumerate}
\end{proof}

\paragraph{The local version.} The same multiplication trick turns the
local form of Brauer's induction theorem into a criterion. For each
prime $\ell$ there is an integer $m$ prime to $\ell$ with $m\cdot1_Q$ a
$\mathbb Z$-combination of inductions from $\ell$-elementary subgroups
(cyclic of order prime to $\ell$, times an $\ell$-group). So
\[
a_{\rho\tau}\in\mathbb Z_{(\ell)}
\iff
\varphi|_E\ \text{is a virtual genuine character for every
$\ell$-elementary }E\le Q.
\]
This twisted version needs no new theory. It follows formally from the
ordinary induction theorems, which answers the hedge in my previous
report.

\paragraph{What remains.} Integrality is now a prime-by-prime problem.
\begin{itemize}
\item \emph{$\ell\ne p$.} The $\ell$-groups map injectively to the finite
metacyclic quotient $H/ZU_D^1$, so the $\ell$-elementary subgroups are
small. They are plausibly controlled by tori and the tame (Adams-type)
arguments. But there is a gap: an abelian subgroup of $Q$ is only
commutative modulo $N$, so it need not lie in one torus. This needs an
argument.
\item \emph{$\ell=p$: the core.} The $p$-elementary subgroups $C\times P$
contain nonabelian $p$-subgroups $P$ of the covered pro-$p$ group
$\widetilde{U_D^1}$. Over that group the cover does not split when
$p\mid n$ (Part~III, Proposition~8.1). This is exactly where no
torus-level information can suffice.
\end{itemize}

\subsection*{3.2 The most promising routes to the $p$-local core}
\begin{enumerate}[label=(\arabic*)]
\item \emph{Explicit construction for $r=2$.} $D$ is quaternion and
$p$ is odd, so $D$ is tamely ramified. Its representations are explicit
(Howe's admissible pairs; Corwin--Howe characters).
\begin{itemize}
\item Build the genuine representations $\tau(E,\tilde\xi)$ by inducing
from $\widetilde{E^\times}\cdot\kappa^*(U_D^m)$, using the
Stone--von Neumann module on the torus.
\item Compute their characters by the Frobenius formula.
\item Match them with $L_D\Theta_\rho$. The note's explicit indicator
form of $R_E$ is exactly what such a comparison produces.
\item By the positivity lemma it is enough to show that each
$L_D\Theta_\rho$ is the character of \emph{some} representation.
\end{itemize}
In my judgment this is now the most realistic path to a complete theorem
for $r=2$.
\item \emph{Kirillov theory on the deep pro-$p$ part.} On deep
congruence subgroups, where the cover splits, the power map is
dilation by $n$, and dilation maps coadjoint orbits to orbits. The
dimension count gives exactly $c_D$ (my previous report, \S3.4). This
targets precisely the $p$-elementary subgroups above.
\item \emph{The global route.} This is the route of Part~IV,
Proposition~8.5, combined with the positivity lemma. It remains
conditional on the split-place and complex-place matching.
\item \emph{Computation.} Compute the $\ell=p$ restriction test, or the
full row, for $(2,3,\mathbb Q_3(\zeta_3))$ at the first few levels. A
failure would be decisive. Success would indicate what the general
construction must look like.
\end{enumerate}

\paragraph{Bottom line.}
\begin{itemize}
\item \emph{$r=2$.} Step~T is finished, and by the corollary of \S3.1
the coefficients are rational. The remaining gap is the $p$-local part
of gluing, which an explicit construction should be able to close.
\item \emph{$r\ge3$.} Step~T itself is still open.
\item \emph{In general.} I would not describe the problem as solvable
``now''. It has been reduced to one clearly identified local question
for $r=2$, and to two questions for $r\ge3$.
\end{itemize}

\section*{4. Suggested next revision of Part~IV}
\begin{enumerate}[label=(\arabic*)]
\item Make the changes from my previous report (Proposition~8.5, $B_d$,
labels).
\item Insert the note's theorem as Proposition~6.6 \textbf{[P]}, with the
consequence for all $\rho$.
\item Add the positivity lemma, with the single-row caveat of \S2.
\item Add the Artin rationality corollary and the $\ell$-local
criterion as \textbf{[P]}, after an independent check.
\item Restate the open problem as: \emph{$p$-local integrality on
$p$-elementary subgroups of the covered pro-$p$ group} (for $r=2$), plus
Step~T (for $r\ge3$).
\end{enumerate}

\end{document}
